% ============================================================
% Section 2: Related Work
% ============================================================

\section{Related Work}
\label{sec:related}

\subsection{Speech Separation}

Neural speech separation has advanced considerably since the introduction
of time-domain models. Conv-TasNet~\cite{luo2019conv} demonstrated that
depthwise separable convolutions with a learned encoder-decoder could
outperform spectrogram-domain methods in Signal-to-Noise Ratio (SNR).
SepFormer~\cite{subakan2021attention} subsequently replaced the convolutional
stack with a hierarchical Transformer, achieving state-of-the-art
Scale-Invariant SNR Improvement (SI-SNRi) on WSJ0-2Mix and WHAM!~\cite{wichern2019wham}.
Multi-channel extensions such as FaSNet~\cite{luo2020end} and
SpatialNet~\cite{wang2023tf} incorporate inter-microphone time differences
to improve separation in reverberant environments—directly applicable to
the four-microphone automotive setup in AISHELL-5.

\subsection{Multi-Speaker ASR}

The Joint CTC/Attention encoder-decoder framework~\cite{watanabe2017hybrid}
pioneered multi-speaker ASR by estimating speaker-attributed transcripts
within a single forward pass. Serialized Output Training (SOT)~\cite{kanda2020serialized}
extended this to handle a variable number of speakers without permutation
ambiguity. More recently, neural diarization approaches such as
EEND~\cite{fujita2019end} and its variants separate the diarization and
recognition tasks via a two-stage pipeline. Whisper~\cite{radford2023robust},
trained on 680k hours of weakly supervised data, achieves near-human WER
on clean speech in 99 languages, making it a natural choice as the ASR
backbone when labeled in-domain data is unavailable.

\subsection{Speaker Diarization and Role Classification}

Speaker diarization assigns a speaker identity to each time segment.
X-vectors~\cite{snyder2018x} and their successor ECAPA-TDNN~\cite{desplanques2020ecapa}
produce compact embeddings that enable clustering-based diarization.
In automotive contexts, role classification (driver vs.\ passenger) carries
additional value for command prioritization: driver commands typically take
precedence over passenger requests. Prior work in automotive ASR has largely
ignored this distinction, treating all occupants equivalently.

\subsection{In-Car Voice Assistants}

Industrial deployments of in-car voice assistants (e.g., Mercedes-Benz
MBUX, BMW Intelligent Personal Assistant) demonstrate commercial interest but
are proprietary and evaluated on internal benchmarks. Academic datasets for
automotive ASR include CHiME-6~\cite{watanabe2020chime6} (dinner party
scenario, four microphones), which shares methodological similarities with
AISHELL-5 but covers a distinct domain. AISHELL-5~\cite{aishell5_2025} is,
to our knowledge, the \emph{first public corpus} specifically designed for
multi-speaker Mandarin ASR in a real vehicle cabin, making it the natural
benchmark for the task we study.

\subsection{Gap Addressed by This Work}

Despite progress in separation and multi-speaker ASR individually,
no prior work has presented a \emph{complete, open-source, benchmarked} system
covering the full chain from raw multi-channel in-car audio to speaker-labeled
intent. Our work closes this gap by coupling pre-trained separation
(SepFormer), pre-trained ASR (Whisper), rule-based role classification, and an
intent engine into a reproducible pipeline evaluated on AISHELL-5.
